\section{zkBridge de Berkeley RDI}
\label{sec:zkbridge_berkeley}

Una cadena de destino no puede consultar directamente el estado de otra cadena, de modo que, antes
de actuar sobre un dato proveniente de ella, necesita establecer que ese dato pertenece a un bloque
aceptado por el consenso de la cadena de origen. Ejecutar en un contrato las reglas de un cliente
ligero evitaría recurrir a un comité externo, pero verificar en destino las firmas y transiciones
que exigen puede resultar prohibitivo. La propuesta zkBridge, de Xie et
al., traslada esa ejecución fuera de la cadena y envía al contrato de destino el nuevo encabezado
junto con una prueba sucinta de que satisface la regla de validación del cliente
ligero~\cite{xie2022zkbridge}.

Los nodos de la red de retransmisión consultan varios nodos completos de la cadena de origen,
obtienen el encabezado que sucede al último registrado y demuestran que la regla de validación,
aplicada al estado vigente del cliente ligero y al encabezado anterior, acepta el nuevo. El contrato
actualizador (\textit{updater contract}) comprueba que el encabezado anterior ya figure en su
registro y que la prueba sea válida respecto del estado que mantiene, y solo entonces actualiza ese
estado e incorpora el encabezado, que las aplicaciones consultan mediante una interfaz independiente
de su lógica. La participación en la retransmisión es abierta: el contrato verifica la prueba sin
necesidad de confiar en quien la presenta. Para que el puente siga incorporando encabezados, debe
haber al menos un retransmisor honesto que los envíe. El análisis de los autores supone además que
la cadena de origen mantiene la consistencia y sigue progresando, que su protocolo de cliente ligero
es seguro y que el sistema de pruebas es sólido~\cite{xie2022zkbridge}.

Aceptar un encabezado no autoriza por sí solo ninguna operación, porque el encabezado únicamente
compromete el contenido del bloque mediante raíces de Merkle. En zkBridge, el contrato receptor de
cada aplicación obtiene el encabezado del actualizador y verifica contra la raíz correspondiente una
prueba adicional de que un dato concreto, como la
inclusión de una transacción, pertenece al bloque. Que un encabezado aceptado no quede luego fuera
de la cadena canónica depende, a su vez, de que el verificador del cliente ligero rechace
encabezados que puedan quedar huérfanos y del criterio de confirmación con que se extrae la cadena
principal del registro~\cite{xie2022zkbridge}.

\begin{sloppypar}
En su tesis de grado en la FP-UNA, Figueredo Camé y Riquelme Godoy reseñaron la propuesta de
Berkeley y construyeron su prototipo a partir de una implementación abierta de
zkBridge~\cite{figueredo2026zkbridge}. El prototipo actual presenta una
separación semejante, en la que \texttt{BeaconLightClient} registra el hash del bloque de ejecución
certificado por la prueba de consenso y \texttt{ZKReceiver} comprueba la inclusión y la ejecución
exitosa de la transacción antes de acuñar (Sección~\ref{sec:puente_arquitectura}). La ruta \lstinline|mintRemote| ilustra esa
división, ya que el receptor aprovecha un anclaje previo sin que se vuelva a presentar la prueba que
lo estableció. Se trata de una semejanza de diseño y no de una equivalencia de garantías, puesto que
el prototipo no encadena cada encabezado con su predecesor, recibe el compromiso del comité de su
propietario y puede anclar bloques no finalizados (Sección~\ref{sec:puente_aceptacion}). La sección siguiente aborda la expresión en Noir de los
enunciados de consenso y de inclusión.
\end{sloppypar}
